% Appendix E: Reproducibility Manifest. Scope-B build chain.

\subsection*{E.1\ Build chain}

The end-to-end reproduction of every figure, table, and numerical claim in this manuscript runs in three phases.

\paragraph{Phase~A --- framework regression.} Verify the \texttt{lcp} reference implementation against \cite{lcp2025}'s worked examples.

\begin{verbatim}
cd workspace
python -m pytest lcp/tests/        # 55 tests, < 1 s
python examples/lcp_demo_v10_2.py  # Example 1 end-to-end
python examples/lcp_demo_example2.py
\end{verbatim}

The regression suite reproduces $w^\dagger = (5.5, 5.5)$, $r^\dagger = 4.5$ for Example~1; $c_1 = 7$, $\kappa_1[3] = 0.75$, witness slack $5.5$ for Example~2. The two demo scripts replay the full pipeline (diagnostics $\to$ Algorithm~1A/1B $\to$ cascade-vs-WS solve $\to$ compliance check $\to$ \S10.2 necessity probe) end-to-end against the paper-published numbers.

\paragraph{Phase~B --- closed-loop dual-condition protocol.} Re-run the 32-scenario protocol on \texttt{nuPlan-mini}.

\begin{verbatim}
cd workspace

# C1 (operational WS-L1; ~2 min/scenario, ~30 min total)
NUPLAN_EXP_ROOT=/workspace/nuplan-project/exp \
  python examples/12_batch_two_level_mpc_planner.py \
    --penalty-form l1 --runtime-mode ws --seed 7 \
    --output-dir examples/outputs/manuscript/C1_instrumented

# C4 (cascade reference; ~14 min/scenario, ~4 h total)
NUPLAN_EXP_ROOT=/workspace/exp_worker_c4/exp \
  python examples/12_batch_two_level_mpc_planner.py \
    --penalty-form l1 --runtime-mode cascade --seed 7 \
    --output-dir examples/outputs/manuscript/C4_cascade
\end{verbatim}

Each batch emits one MP4 + summary PNG + per-tick \texttt{*\_log.csv} per scenario, plus a per-batch \texttt{batch\_summary\_l1.csv}.

\paragraph{Phase~C --- artefact generation + LaTeX build.}

\begin{verbatim}
cd workspace
bash examples/regen_manuscript.sh   # diagnostics + necessity scans
                                    # + figures + tables to IEEE_T-IV/

cd ../IEEE_T-IV
pdflatex Main && bibtex Main \
  && pdflatex Main && pdflatex Main
\end{verbatim}

The \texttt{regen\_manuscript.sh} script invokes \texttt{scan\_diagnostics.py}, \texttt{scan\_necessity.py}, and \texttt{make\_manuscript\_artefacts.py}, which mirror the generated figures into \texttt{IEEE\_T-IV/Figures/} and the \texttt{\textbackslash input\{\}}able tables into \texttt{IEEE\_T-IV/Sections/tables/}.

\subsection*{E.2\ Software versions}

The deployment uses the Python scientific stack with the following pinned versions (the project \texttt{requirements.txt} is the authoritative source).

\begin{table}[h]
    \centering
    \footnotesize
    \begin{tabularx}{\linewidth}{@{}l l X@{}}
        \toprule
        Package & Version & Role \\
        \midrule
        Python & 3.9.23 & runtime \\
        numpy / scipy & 1.x / 1.7+ & linear algebra, \texttt{scipy.optimize.linprog} (HiGHS) for Algorithms 0/1A/1B \\
        casadi & 3.5+ & nonlinear-program modelling and IPOPT interface \\
        ipopt & 3.14+ & interior-point inner-OCP solver \\
        nuplan-devkit & per pin & closed-loop simulator, scenario filter, \texttt{run\_simulation.py} \\
        matplotlib & 3.x & figure rendering (\texttt{scripts/ieee\_style.py} for IEEEtran column widths) \\
        pytest, hypothesis & latest & regression test suite \\
        \bottomrule
    \end{tabularx}
\end{table}

\subsection*{E.3\ Hardware spec}

The protocol was executed on a single workstation. Each batch runs \texttt{worker=sequential} (one scenario at a time per process) and the two batches can be parallelised across distinct \texttt{NUPLAN\_EXP\_ROOT} directories. Memory peak per process is dominated by the nuPlan SimulationLog buffer ($\sim 1$~GB per 15-second scenario). Disk footprint per scenario is $\sim 50$~MB (MP4 + summary PNG + per-tick CSV); the full 32-scenario protocol fits in a few GB.
